% ============================================
% Johns Hopkins Letter Template
% Assistant Teaching Professor Cover Letter – William & Mary
% ============================================

\documentclass[11pt,letterpaper]{letter}

\usepackage{graphicx}
\usepackage{eso-pic}
\usepackage{geometry}
\usepackage{microtype}
\usepackage[hidelinks]{hyperref}

% ----- Page Setup -----
\geometry{left=25mm,right=25mm,top=25mm,bottom=25mm}

% ----- Logo Placement (foreground, first page only) -----
\newcommand{\PutJHULogoFG}[1]{%
  \AddToShipoutPictureFG*{%
    \AtPageUpperLeft{%
      \hspace{10mm}%
      \raisebox{-38mm}{%
        \includegraphics[width=6cm]{#1}%
      }%
    }%
  }%
}

% ----- Sender Information -----
\signature{Decory Edwards}
\address{Department of Economics \\ Johns Hopkins University \\ Baltimore, MD 21218 \\ \texttt{dedwar65@jh.edu}}

\begin{document}
\PutJHULogoFG{Images/jhulogo.png}

\begin{letter}{Search Committee \\
Department of Economics \\
College of William \& Mary \\
Chancellors Hall \\
Williamsburg, VA 23187 \\ USA}

\opening{Dear Members of the Search Committee,}

I am writing to express my interest in the Assistant Teaching Professor of Economics position in the Department of Economics at William \& Mary. I am a Ph.D. candidate in Economics at Johns Hopkins University, advised by Professors Christopher D. Carroll, Nicholas W. Papageorge, and Stelios Fourakis, and I expect to receive my degree in May 2026. My research fields are macroeconomics, wealth and income dynamics, and computational economics.

My research agenda focuses on understanding the determinants of wealth inequality and the mechanisms through which heterogeneity in returns to wealth shapes aggregate outcomes. In my job-market paper, I estimate the degree of heterogeneity in individual portfolio returns using micro-data from the Survey of Consumer Finances and simulate a structural model in which households face idiosyncratic return risk. I show that allowing for persistent heterogeneity in returns substantially improves the model’s ability to match the empirical distribution of wealth, offering a tractable way to connect micro-level return dispersion to macroeconomic inequality.

In a second project, I use the Health and Retirement Study to examine the relationship between interpersonal trust and portfolio performance. Building on the literature that finds a hump-shaped relationship between trust and income, I show that a similar relationship holds for individual returns to wealth measured in the HRS data, highlighting a behavioral channel through which social attitudes shape saving and investment outcomes. This work also provides new estimates of the persistent component of returns to net worth, which motivate the calibration of heterogeneity in my structural model.

William \& Mary’s Department of Economics has invited applications for a three-year instructional appointment with a 3-3 teaching load, an emphasis on exceptional teaching in macroeconomics (introductory and intermediate), and service responsibilities commensurate with rank. I am excited about the opportunity to join a department with a strong commitment to undergraduate education, large-section teaching, and a framework for teaching-faculty promotion. My teaching-centered research profile and experience in quantitative methods position me well to contribute to the introduction and intermediate macroeconomics courses, as well as to support the department’s broader instructional mission.

\textbf{Teaching.} My teaching experience centers on undergraduate macroeconomics and an upper-level macro-policy seminar, where I have led weekly discussion sections and held regular office hours for classes ranging from 16 to 40 students. My approach is student-centered: I aim to help students “convince themselves” that they have moved from confusion to understanding by holding structured office-hour coaching that builds trust and confidence. At William \& Mary, I am prepared to teach \textit{Principles of Macroeconomics}, \textit{Intermediate Macroeconomic Theory}, \textit{Money and Banking}, and \textit{Quantitative Methods for Economics}. I would also welcome developing electives such as \textit{Economics of Inequality and Tax Policy} or \textit{Computational Macroeconomics and Policy}, emphasizing reproducible workflows and evidence-based decision-making.

I am committed to inclusive pedagogy and student success in line with William \& Mary’s mission as a public university dedicated to excellence and inclusive undergraduate education. In my courses, I build multiple modes of assessment, scaffold quantitative skills for diverse learners, and mentor students in research and applied data projects to ensure all students can thrive.

I would be honored to join the Department of Economics at William \& Mary and to contribute to its teaching mission, service responsibilities, and student-development goals. Thank you for your consideration; I welcome the opportunity to discuss my fit with the department at your convenience.

\closing{Sincerely,}

\end{letter}
\end{document}
